\section{Ассоциативные контейнеры \texttt{std::set} и \texttt{std::map}}

\subsection{Какие основные характеристики имеют ассоциативные контейнеры?}
Ассоциативные контейнеры -- это контейнеры, в которых элементы хранятся не в порядке вставки, а в соответствии с некоторым правилом упорядочивания. Для стандартных контейнеров \texttt{std::set} и \texttt{std::map} таким правилом обычно является сравнение по \texttt{operator<} (или по пользовательскому компаратору).

\begin{definition}[Ассоциативный контейнер]
Ассоциативный контейнер хранит элементы в упорядоченном виде и поддерживает эффективный поиск по ключу. Основные операции -- поиск, вставка и удаление -- обычно имеют логарифмическую сложность.
\end{definition}

Ключевые свойства:
\begin{itemize}
    \item элементы упорядочены по строгому слабому порядку;
    \item поиск выполняется существенно быстрее, чем линейный просмотр;
    \item контейнер не является последовательным: порядок обхода задаётся сортировкой, а не временем вставки;
    \item обычно не допускаются дубликаты ключей;
    \item итерация по контейнеру идёт в отсортированном порядке.
\end{itemize}

\subsection{Что такое \texttt{std::set}? Каковы основные операции и ограничения?}
\texttt{std::set} -- это множество уникальных значений. Один и тот же элемент не может присутствовать в контейнере дважды.

\begin{definition}[\texttt{std::set}]
\texttt{std::set<T>} -- ассоциативный контейнер, хранящий элементы типа \texttt{T} в упорядоченном виде без дубликатов.
\end{definition}

Основные операции:
\begin{itemize}
    \item \texttt{find(x)} -- поиск элемента;
    \item \texttt{insert(x)} -- вставка элемента;
    \item \texttt{erase(x)} или \texttt{erase(it)} -- удаление элемента;
    \item \texttt{begin()}, \texttt{end()} -- обход контейнера.
\end{itemize}

Если вставляется значение, уже присутствующее в \texttt{set}, то новый элемент не добавляется: контейнер остаётся без изменений.

\begin{remark}
В учебной формулировке часто говорят, что элементы \texttt{set} должны быть \textit{default constructable}, \textit{copyable} и \textit{comparable}. В более точном стандарте важно, чтобы элементы можно было сравнивать с помощью компаратора, задающего \textit{strict weak ordering}. На практике это обычно означает наличие корректного \texttt{operator<} или пользовательского сравнения.
\end{remark}

\subsection{Какие требования накладываются на тип элементов \texttt{std::set} и почему?}
Ключевое требование: элементы должны быть сравнимы через компаратор контейнера так, чтобы задавался строгий слабый порядок.

Это необходимо, потому что \texttt{set} хранит элементы в упорядоченной древовидной структуре и использует сравнение для:
\begin{itemize}
    \item поиска пути в дереве;
    \item проверки эквивалентности ключей;
    \item поддержания корректного порядка при вставке и удалении.
\end{itemize}

На практике элемент обычно должен быть копируемым или перемещаемым и корректно разрушаемым, чтобы контейнер мог безопасно управлять временем жизни узлов.

\subsection{Как итерироваться по \texttt{std::set}? К какой категории относятся итераторы?}
Итерирование по \texttt{std::set} происходит в отсортированном порядке, то есть в том порядке, в каком элементы лежат в дереве поиска.

Пример:
\begin{verbatim}
std::set<int> s = {3, 1, 4};
for (auto it = s.begin(); it != s.end(); ++it) {
    std::cout << *it << ' ';
}
\end{verbatim}

Итераторы \texttt{std::set} относятся к категории \textbf{bidirectional iterator}. Это означает, что:
\begin{itemize}
    \item можно двигаться вперёд через \texttt{++it};
    \item можно двигаться назад через \texttt{--it};
    \item нельзя делать произвольный доступ, как у random access iterator, то есть нельзя писать \texttt{it + n}.
\end{itemize}

Причина в том, что \texttt{set} обычно реализован как дерево, а не как массив.

\subsection{Как искать элемент в \texttt{std::set}? Что возвращает \texttt{find()}?}
Основная операция поиска -- \texttt{find(x)}.

Она возвращает итератор:
\begin{itemize}
    \item на найденный элемент, если он есть в контейнере;
    \item на \texttt{end()}, если элемента нет.
\end{itemize}

Это значение можно использовать:
\begin{itemize}
    \item для проверки наличия элемента;
    \item для удаления через \texttt{erase(it)};
    \item для доступа к найденному элементу через разыменование итератора.
\end{itemize}

Идиома проверки наличия:
\begin{verbatim}
if (s.find(x) != s.end()) {
    // x есть в множестве
}
\end{verbatim}

\subsection{Как работают \texttt{insert()} и \texttt{erase()} в \texttt{std::set}? Почему их сложность связана с \texttt{find()}?}
Обе операции опираются на поиск места в дереве.

\begin{itemize}
    \item \texttt{insert(x)} сначала находит позицию, куда должен попасть элемент, а затем вставляет его.
    \item \texttt{erase(x)} сначала находит нужный элемент, а затем удаляет его и восстанавливает свойства дерева.
\end{itemize}

Поэтому алгоритмическая сложность этих операций определяется сложностью поиска:
\[
O(\log n)
\]
при типичной реализации на сбалансированном дереве.

Если удаление или вставка были бы возможны без поиска, это всё равно не отменяло бы необходимость определить место элемента в упорядоченной структуре. Именно поэтому \texttt{find()} является базовой операцией, от которой зависят остальные.

\subsection{Как обычно реализован \texttt{std::set}?}
Обычная реализация -- это \textbf{сбалансированное бинарное дерево поиска}, чаще всего красно-чёрное дерево (\textit{red-black tree}).

Следствия такой реализации:
\begin{itemize}
    \item поиск, вставка и удаление имеют сложность \texttt{O(log n)};
    \item обход контейнера идёт в отсортированном порядке;
    \item память расходуется на узлы дерева, а не на плотный массив.
\end{itemize}

\begin{remark}
Именно дерево, а не массив, объясняет тип итераторов: можно идти к следующему и предыдущему элементу, но нельзя обращаться к элементу по индексу.
\end{remark}

\subsection{Что такое \texttt{std::pair}? Почему он нужен?}
\texttt{std::pair} -- это вспомогательный класс для хранения двух значений, обычно разных типов.

\begin{definition}[\texttt{std::pair}]
\texttt{std::pair<T1, T2>} -- агрегат, представляющий упорядоченную двойку значений \texttt{first} и \texttt{second}.
\end{definition}

Его используют, когда нужно хранить связанные данные:
\begin{itemize}
    \item ключ и значение;
    \item координаты;
    \item две связанные величины.
\end{itemize}

Свойства \texttt{pair}:
\begin{itemize}
    \item это кортеж из двух элементов;
    \item удобно создаётся через фигурные скобки или \texttt{std::make\_pair};
    \item сравнение идёт лексикографически: сначала \texttt{first}, затем \texttt{second}.
\end{itemize}

\begin{example}
\begin{verbatim}
std::pair<int, std::string> p{1, "one"};
std::cout << p.first << ' ' << p.second;
\end{verbatim}
\end{example}

\subsection{Что такое \texttt{std::map}? Чем он отличается от \texttt{std::set<std::pair<K,V>>}?}
\texttt{std::map} -- это ассоциативный контейнер, который хранит пары \texttt{(key, value)}. Ключи уникальны.

\begin{definition}[\texttt{std::map}]
\texttt{std::map<K, V>} -- контейнер, в котором каждому ключу типа \texttt{K} соответствует значение типа \texttt{V}. Ключи упорядочены, а для каждого ключа хранится не более одного значения.
\end{definition}

Различие между \texttt{std::map<K, V>} и \texttt{std::set<std::pair<K,V>>}:
\begin{itemize}
    \item в \texttt{map} ключ и значение концептуально разделены на \texttt{key} и \texttt{mapped value};
    \item у \texttt{set<pair<...>>} сортировка идёт по всей паре целиком, то есть по \texttt{first}, а при равенстве -- по \texttt{second};
    \item \texttt{map} предоставляет удобный доступ по ключу через \texttt{operator[]};
    \item \texttt{map} моделирует ассоциативный массив, а \texttt{set<pair<...>>} -- просто множество пар.
\end{itemize}

\subsection{Как \texttt{std::map} моделирует ассоциативный массив?}
\texttt{std::map<K, V>} можно рассматривать как частично определённую дискретную функцию
\[
m : K \mapsto V.
\]
То есть не каждому ключу обязательно сопоставлено значение, но для существующих ключей соответствие определено.

Основные операции такие же, как у \texttt{set}:
\begin{itemize}
    \item \texttt{find(key)} -- поиск по ключу;
    \item \texttt{insert(...)} -- вставка пары;
    \item \texttt{erase(...)} -- удаление по ключу или итератору.
\end{itemize}

\subsection{Как работает \texttt{operator[]} у \texttt{std::map}? Почему важна неконстантная перегрузка?}
Операция \texttt{m[k]} возвращает ссылку на значение, связанное с ключом \texttt{k}.

Если ключ уже есть, возвращается ссылка на его значение.  
Если ключа нет, то в контейнер автоматически вставляется новая пара \texttt{(k, V())}, где \texttt{V()} -- значение, созданное по умолчанию.

\begin{example}
\begin{verbatim}
std::map<int, int> m;
int a = m[42];   // a == 0, а в map появится ключ 42 со значением 0
\end{verbatim}
\end{example}

\begin{remark}
В стандартном \texttt{std::map} оператор \texttt{operator[]} является только \textbf{неконстантным}. Константной перегрузки нет, потому что доступ через \texttt{[]} может изменить контейнер, вставив новый элемент. Для чтения из \texttt{const std::map} обычно используют \texttt{find()} или \texttt{at()}.
\end{remark}

Поэтому:
\begin{itemize}
    \item для \texttt{std::map<K, V> m} выражение \texttt{m[k]} допустимо;
    \item для \texttt{const std::map<K, V> m} выражение \texttt{m[k]} недопустимо (ошибка компиляции).
\end{itemize}

Смысл \texttt{operator[]} -- сделать \texttt{map} похожим на массив: можно писать \texttt{m[k] = v}. Это удобно, но нужно помнить, что при отсутствии ключа происходит вставка.

\subsection{Как итерироваться по \texttt{std::map}? Что получается при разыменовании итератора?}
Итераторы \texttt{std::map} тоже относятся к категории \textbf{bidirectional iterator}. То есть по ним можно двигаться вперёд и назад, но нельзя прыгать на произвольное расстояние.

При разыменовании итератора получается объект типа \texttt{std::pair<const K, V>}:
\begin{itemize}
    \item \texttt{first} -- ключ;
    \item \texttt{second} -- значение.
\end{itemize}

Ключ обычно объявлен как \texttt{const}, чтобы нельзя было изменить порядок элементов в контейнере через итератор.

\subsection{Что происходит при вставке элемента в \texttt{std::map}, если ключ уже существует?}
Тут поведение зависит от способа вставки.

\begin{itemize}
    \item \texttt{m[k] = v} -- если ключ уже есть, то изменится только связанное значение;
    \item \texttt{m.insert(\{k, v\})} -- если ключ уже есть, новая пара не будет вставлена, и существующее значение не изменится.
\end{itemize}

То есть:
\begin{itemize}
    \item \texttt{operator[]} подходит для обновления значения по ключу;
    \item \texttt{insert()} подходит для вставки только тогда, когда ключ отсутствует.
\end{itemize}

\subsection{Как правильно искать элемент в \texttt{std::map}?}
Идиоматический способ -- использовать \texttt{find(key)}.

\begin{verbatim}
auto it = m.find(k);
if (it != m.end()) {
    // найдено: it->first == k, it->second -- значение
}
\end{verbatim}

Такой подход полезен, потому что:
\begin{itemize}
    \item не создаёт новый элемент;
    \item работает без лишней вставки;
    \item сразу даёт доступ к найденной паре.
\end{itemize}

\begin{remark}
Не стоит использовать \texttt{m[k]} для проверки наличия ключа, если не нужна вставка. Иначе при отсутствии ключа в контейнере появится новый элемент с значением по умолчанию.
\end{remark}

\subsection{Чем отличаются \texttt{std::multiset} и \texttt{std::multimap} от \texttt{std::set} и \texttt{std::map}?}
\texttt{std::multiset} и \texttt{std::multimap} устроены так же, как соответствующие упорядоченные контейнеры, но разрешают дубликаты ключей.

\begin{itemize}
    \item \texttt{std::multiset} хранит несколько одинаковых значений;
    \item \texttt{std::multimap} хранит несколько значений под одним и тем же ключом;
    \item для поиска таких элементов часто используют \texttt{equal\_range()}.
\end{itemize}

\begin{remark}
Если в задаче возможны повторяющиеся ключи, нужен не \texttt{set}/\texttt{map}, а \texttt{multiset}/\texttt{multimap}.
\end{remark}

\subsection{Чем упорядоченные ассоциативные контейнеры отличаются от неупорядоченных?}
Упорядоченные контейнеры хранят элементы в порядке компаратора. Неупорядоченные контейнеры используют хеш-таблицу.

\begin{itemize}
    \item \texttt{std::unordered\_set} и \texttt{std::unordered\_map} не поддерживают сортированный обход;
    \item поиск в среднем близок к \texttt{O(1)};
    \item порядок обхода зависит от хеш-таблицы и не обязан быть стабильным;
    \item при росте числа элементов возможен \texttt{rehash}, который перераспределяет элементы по корзинам.
\end{itemize}

\begin{remark}
В неупорядоченных контейнерах вместо компаратора используются хеш-функция и проверка равенства ключей. Коллизии допустимы и обрабатываются внутри таблицы.
\end{remark}

\subsection{Какая типичная внутренняя реализация у \texttt{std::map} и какова асимптотика операций?}
Как и \texttt{std::set}, \texttt{std::map} обычно реализован как сбалансированное бинарное дерево поиска, чаще всего красно-чёрное дерево.

Отсюда следуют асимптотики:
\begin{itemize}
    \item поиск -- \texttt{O(log n)};
    \item вставка -- \texttt{O(log n)};
    \item удаление -- \texttt{O(log n)};
    \item обход всех элементов -- \texttt{O(n)}.
\end{itemize}

Именно эта реализация делает \texttt{map} удобным ассоциативным массивом: быстрый поиск по ключу и упорядоченный обход по возрастанию ключей.
